%%
%% CSD Assignment 1 - Final Technical Report
%% Based on the ACM "acmart" manuscript (single-column) style.
%% Course: Computer System Design (CSD) | Instructor: Shivam Kushwaha
%% Student: Katari Venu (Roll: 12341110) | IIT Bhilai
%%
\documentclass[manuscript,screen,nonacm]{acmart}

\usepackage{booktabs}
\usepackage{tabularx}
\usepackage{subcaption}
\usepackage{listings}
\usepackage{xcolor}

\AtBeginDocument{%
  \providecommand\BibTeX{{Bib\TeX}}}

\begin{document}

%% =================== TITLE & AUTHOR BLOCK =========================
\title{AI-based Village Pond Planning System}
\subtitle{CSD Assignment 1 -- Final Technical Report}

\author{Katari Venu}
\affiliation{%
  \institution{Department of Computer Science and Engineering, Indian Institute of Technology Bhilai}
  \city{Durg}
  \state{Chhattisgarh}
  \country{India}}
\email{katariv@iitbhilai.ac.in}

\renewcommand{\shortauthors}{Katari Venu}

%% =================== ABSTRACT ======================================
\begin{abstract}
Water conservation in rural agricultural watersheds is an urgent priority across peninsular and central India, where irregular monsoon precipitation causes severe seasonal water distress. Decentralized farm ponds offer an effective, low-cost rainwater harvesting solution; however, manual site selection is hampered by undulating terrain, complex drainage topologies, and a lack of integrated hydrometeorological modeling. This report presents the design, implementation, and evaluation of an end-to-end, AI- and geospatial-assisted web application for automated village pond planning. The system combines digital elevation modeling (DEM), deterministic eight-node (D8) steepest-descent flow direction routing, automated topological upstream catchment delineation, and morphological valley-floor hazard exclusion with real-time and climatological precipitation data from Open-Meteo and the India Meteorological Department (IMD). 

Through an interactive GIS frontend built on Leaflet.js, village administrators can select arbitrary land parcels on high-resolution satellite imagery or upload KML/KMZ contour maps. The system dynamically computes the optimal pond location, delineates the contributing catchment area, calculates harvestable water volume via the Rational Method, and recommends civil engineering pond dimensions conforming to Indian Standards (IS: 4987). Tested on a real 1-meter contour dataset from the Shivnath River watershed near Bhilai (159,113 terrain points across 1,355 contour lines), the system identified an optimal interior pond site at $21.2488^\circ\text{ N}, 81.3004^\circ\text{ E}$ with a contributing catchment of $63,001.02\text{ m}^2$ ($6.30\text{ ha}$), harvestable monsoon runoff volume of $22,683\text{ m}^3$ ($22.68\text{ ML}$), and recommended a trapezoidal storage reservoir of $6,804\text{ m}^3$ capable of irrigating $7.5\text{ ha}$ of rabi crops. Deployed over the IIT Bhilai campus network infrastructure across mapped application ports, the backend delivers sub-200 millisecond response times under active caching.
\end{abstract}

\keywords{Village Pond Planning, Geospatial Hydrology, Rainwater Harvesting, Catchment Delineation, D8 Flow Routing, Web GIS, REST APIs, Computer System Design}

\maketitle

%% =====================================================================
\section{Introduction}
\label{sec:intro}
Water scarcity during post-monsoon months poses a critical socioeconomic bottleneck to agricultural productivity in rural India. In rain-fed farming ecosystems, over 80\% of annual precipitation occurs within the four-month southwest monsoon (June to September). Without adequate local capture mechanisms, the vast majority of surface runoff escapes into regional river corridors or causes soil erosion, leaving rural communities vulnerable to crop failure and acute groundwater table depletion during the rabi season. Under national initiatives such as the Jal Shakti Abhiyan and Mission Amrit Sarovar, the construction of village ponds (farm ponds and community percolation tanks) has emerged as the foremost sustainable intervention to store excess runoff and recharge local aquifers.

Despite their benefits, the actual efficacy of village ponds depends entirely on spatial placement. Constructing a pond in an arbitrary depression often leads to structural disaster: if placed directly inside an active river or torrent bed, the embankment is washed away during peak flood events; conversely, if placed on an elevated ridge or outside natural drainage pathways, the pond remains dry due to an insufficient contributing catchment area. The objective of this assignment is to engineer an automated, full-stack, AI- and geospatial-assisted web application that ingests contour datasets or user-defined map boundaries, algorithmically determines optimal, safe pond excavation sites, delineates upstream drainage basins, models harvestable rainfall runoff, and provides actionable civil engineering dimensions overlaid on interactive satellite maps.

The remainder of this report is organized as follows: Section~\ref{sec:requirements} defines the functional and non-functional requirements. Section~\ref{sec:architecture} details the system architecture and technology choices. Section~\ref{sec:methodology} describes the mathematical and algorithmic methodology for terrain processing, D8 routing, rainfall integration, and runoff estimation. Section~\ref{sec:implementation} documents backend API endpoints, frontend GIS visualization, and caching strategies. Section~\ref{sec:csd-themes} explicitly maps our implementation to core Computer System Design (CSD) themes. Section~\ref{sec:results} presents experimental results and performance benchmarks, followed by limitations in Section~\ref{sec:discussion}, AI usage in Section~\ref{sec:ai-usage}, and deployment source code access in the Appendix.

\subsection{Motivation}
Manual site selection for village ponds is historically slow, error-prone, and expensive. Field survey teams typically rely on visual inspection or rudimentary handheld GPS surveys without access to comprehensive elevation grids or flow routing algorithms. In rural administrations, three fundamental challenges make manual selection unfeasible at scale:
\begin{enumerate}
    \item \textbf{Terrain and Hydrological Complexity}: Surface runoff does not follow human administrative borders or linear slopes. Accurately determining which ridges and micro-depressions feed a specific excavation point requires raster-level flow accumulation that cannot be performed by eye.
    \item \textbf{Decoupled Meteorological Information}: Village planners rarely have centralized, location-specific precipitation time series to evaluate whether the upstream catchment area will yield sufficient harvestable volume to justify the excavation cost.
    \item \textbf{Environmental and Disaster Risks}: Placing ponds in the lowest terrain points frequently places them within active flash-flood channels or river floodplains, risking civil embankment failure. An automated system can enforce mathematical buffer zones that exclude hazardous floodways while retaining proximity to high-yield runoff pathways.
\end{enumerate}

\subsection{Scope of the Project}
The developed system covers:
\begin{itemize}
    \item Interactive web GIS with high-resolution satellite imagery (Esri World Imagery) and cartographic layers.
    \item Flexible land area selection via interactive on-map bounding box drawing, village presets, or direct KML/KMZ contour map upload.
    \item Automated Digital Elevation Model (DEM) construction via metric coordinate projection, Delaunay triangulation, nearest-neighbor boundary filling, and Gaussian smoothing.
    \item Deterministic Eight-Neighbor (D8) downhill flow direction calculation and upstream catchment area accumulation.
    \item Automated morphological exclusion of natural valley floors and riverbed corridors to avoid structural flood risks.
    \item Multi-criteria optimization scoring (upstream area, elevation relative to terrain, and local slope) to select the optimal pond location.
    \item Upstream contributing drainage basin delineation and GeoJSON polygon generation.
    \item Integration of historical annual, seasonal (monsoon), and monthly precipitation data via the Open-Meteo Historical Archive API with verified India Meteorological Department (IMD) climatological normal fallbacks.
    \item Surface runoff estimation using the Rational Method ($Q = C \cdot P \cdot A$) and civil dimensioning (depth, surface area, side slopes, excavation volume, and irrigation potential) conforming to Indian Standard IS: 4987.
    \item Interactive visual overlay of the suggested pond site, catchment polygon, land boundary, exclusion buffer, and monthly runoff analytics.
\end{itemize}

\noindent \textbf{Out of Scope}: The system does not conduct geotechnical core-drilling soil permeability tests, does not perform structural finite-element concrete reinforcement design for spillways, and does not interface with municipal cadastral property deed registries.

%% =====================================================================
\section{Problem Statement and Requirements}
\label{sec:requirements}
The problem statement requires engineering an end-to-end web system that enables village administrators to assess land for rainwater harvesting. Table~\ref{tab:requirements} maps each mandatory functional requirement to the concrete module and source file in our codebase.

\begin{table}[h]
  \caption{Functional requirements and where they are implemented}
  \label{tab:requirements}
  \begin{tabular}{p{5.5cm}p{7.5cm}}
    \toprule
    Requirement & Implemented in Module / File \\
    \midrule
    Satellite imagery display & \texttt{frontend/index.html} (Leaflet.js + Esri World Imagery Tile Layer) \\
    Contour map visualization & \texttt{main.py} (\texttt{parse\_contours}, \texttt{generate\_sample\_contours\_geojson}) \\
    Available-land identification & \texttt{main.py} (\texttt{create\_valley\_exclusion\_mask}, slope filtering) \\
    Catchment area estimation & \texttt{main.py} (\texttt{calculate\_flow\_direction}, \texttt{calculate\_upstream\_area}) \\
    Historical rainfall query & \texttt{main.py} (\texttt{get\_historical\_rainfall}, Open-Meteo \& IMD Normals) \\
    Runoff volume estimation & \texttt{main.py} (\texttt{calculate\_runoff\_and\_pond\_sizing}, Rational Method) \\
    Pond depth / storage recommendation & \texttt{main.py} (\texttt{calculate\_runoff\_and\_pond\_sizing}, IS: 4987 Civil Engine) \\
    Combined overlay / results view & \texttt{frontend/index.html} (GeoJSON layers, HUD Dashboard, Chart.js) \\
    \bottomrule
  \end{tabular}
\end{table}

\subsection{Non-Functional Requirements}
\begin{enumerate}
    \item \textbf{Performance and Latency}: Dynamic land area selection queries on the map must execute in under $300\text{ ms}$. Complex KML files containing over 150,000 coordinate points must be parsed, interpolated into a DEM, and processed through D8 routing in under $6\text{ seconds}$.
    \item \textbf{Scalability and Concurrency}: The backend must utilize an asynchronous event loop (ASGI) capable of handling concurrent client map interactions without blocking computational threads.
    \item \textbf{Infrastructure and Port Mapping Constraints}: The solution must operate strictly within the allotted host environment at IP \texttt{10.1.75.79} across designated port families ($2000, 3000, 4000, 5000, 6000, 7000$) mapped to public campus access ports $X289$ (e.g., port $6000 \to 6289$, port $3000 \to 3289$).
    \item \textbf{High Availability and Fault Tolerance}: The backend must remain operational 24/7 as background daemons with automated restart capabilities. If external weather APIs face rate limits or network outages, the system must degrade gracefully to verified regional climatological data.
    \item \textbf{Usability and Cross-Platform Accessibility}: The frontend must be a zero-install, responsive Single Page Application (SPA) compatible with modern mobile, tablet, and desktop web browsers.
\end{enumerate}

%% =====================================================================
\section{System Architecture and High-Level Design}
\label{sec:architecture}
The application adopts a decoupled, service-oriented client-server architecture designed for high responsiveness, maintainability, and network flexibility. Figure~\ref{fig:architecture} illustrates the high-level architecture and data flow across network tiers.

\begin{figure}[h]
  \centering
  \begin{verbatim}
  +--------------------------------------------------------------------------+
  |                        CLIENT TIER (Web Browser)                         |
  |  Leaflet.js Map  |  Esri Satellite Layer  |  Results HUD  |  Chart.js    |
  +--------------------------------------------------------------------------+
           |                                                      ^
    HTTP Request (Bounds / KML)                        GeoJSON Overlays & KPIs
           |                                                      |
           v                                                      |
  +--------------------------------------------------------------------------+
  |              CAMPUS NETWORK GATEWAY (Port Forwarding X000 -> X289)       |
  |  Frontend URL: http://10.1.75.79:6289   <--->   Host Port 6000           |
  |  Backend API:  http://10.1.75.79:3289   <--->   Host Port 3000           |
  +--------------------------------------------------------------------------+
           |                                                      ^
           v                                                      |
  +--------------------------------------------------------------------------+
  |                    BACKEND TIER (FastAPI / Uvicorn ASGI)                 |
  |  [Router] /api/analyze-area  |  /analyzeContour  |  /api/rainfall        |
  |  [Validation] Pydantic Data Models  |  CORS Middleware                   |
  +--------------------------------------------------------------------------+
           |                                                      ^
           v                                                      |
  +--------------------------------------------------------------------------+
  |                   HYDROLOGICAL COMPUTATION ENGINE                        |
  |  Metric Projection  |  Interpolated DEM (SciPy)  |  D8 Flow Routing      |
  |  Upstream Accumulation  |  Valley Hazard Exclusion Mask                  |
  |  Convex Hull Catchment Delineation  |  Rational Runoff Modeling (IS 4987)|
  +--------------------------------------------------------------------------+
           |                                                      ^
           v                                                      |
  +--------------------------------------------------------------------------+
  |              PERSISTENCE & EXTERNAL DATA INTEGRATION                     |
  |  In-Memory LRU Cache  |  IMD Normals  |  Open-Meteo Historical Archive   |
  +--------------------------------------------------------------------------+
  \end{verbatim}
  \caption{High-level system architecture and network data flow.}
  \label{fig:architecture}
  \Description{Block diagram showing the flow from Client Tier to Network Gateway to Backend FastAPI to Hydrological Computation Engine and External Data Providers.}
\end{figure}

\subsection{Technology Stack}
The system leverages a specialized scientific and geospatial technology stack:
\begin{itemize}
    \item \textbf{Backend Web Framework}: \textbf{FastAPI (Python 3.12)} running on the \textbf{Uvicorn} ASGI server. FastAPI was chosen over Flask due to its native asynchronous execution model, automatic OpenAPI (Swagger) generation at \texttt{/docs}, strict Pydantic type validation, and high performance for concurrent API clients.
    \item \textbf{Scientific and Spatial Computing}: \textbf{NumPy} for vectorized multidimensional matrix operations, \textbf{SciPy (interpolate, ndimage, spatial)} for 2D Delaunay grid interpolation, Gaussian smoothing filters, morphological binary dilation for hazard buffers, and \texttt{ConvexHull} for polygon boundary extraction.
    \item \textbf{Frontend GIS Framework}: \textbf{HTML5, Vanilla CSS3, and JavaScript (ES6+)} utilizing \textbf{Leaflet.js (v1.9.4)} and \textbf{Leaflet-Draw} for zero-dependency interactive mapping. Tile layers are streamed directly from \textbf{Esri World Imagery} and OpenStreetMap.
    \item \textbf{Visualization and Telemetry}: \textbf{Chart.js (v4.4)} for rendering interactive monthly rainfall versus harvestable runoff bar-and-line graphs.
    \item \textbf{External Climatology}: \textbf{Open-Meteo Historical Weather Archive REST API} integrated with asynchronous fallback to official \textbf{India Meteorological Department (IMD)} climatological normals for Chhattisgarh.
\end{itemize}

%% =====================================================================
\section{Methodology}
\label{sec:methodology}
This section details the mathematical algorithms and data pipelines implemented in the hydrological processing engine.

\subsection{Terrain and Elevation Analysis}
Terrain modeling begins with discrete elevation coordinates extracted from uploaded KML/KMZ files or geographic bounding boxes. Because raw spatial coordinates are expressed in angular WGS84 degrees (longitude $\lambda$, latitude $\phi$), calculating metric slopes and distances directly causes severe latitude-dependent distortions. To eliminate distortion, coordinates are projected into a local Cartesian metric coordinate frame $(x, y)$ centered at origin $(\lambda_0, \phi_0)$:
\begin{equation}
x = (\lambda - \lambda_0) \cdot 111320.0 \cdot \cos\left(\frac{\phi_0 \cdot \pi}{180}\right), \quad y = (\phi - \phi_0) \cdot 110540.0
\end{equation}
where $111320.0\text{ m/deg}$ represents the equatorial longitude scale and $110540.0\text{ m/deg}$ represents the meridional latitude scale. Duplicate spatial points are averaged into unique $(x, y)$ coordinates.

A continuous Digital Elevation Model (DEM) raster grid $Z(r, c)$ of dimensions $N \times N$ ($N=180$) is constructed across the bounding interval $[x_{\min}, x_{\max}] \times [y_{\min}, y_{\max}]$. Interpolation is performed using SciPy's \texttt{griddata} via 2D Delaunay triangulation (linear interpolation). Boundary extrapolations and edge gaps are resolved using nearest-neighbor interpolation. To suppress high-frequency digitization artifacts and measurement noise while preserving regional terrain relief, a Gaussian smoothing filter is applied:
\begin{equation}
Z_{\text{smooth}} = Z * G_\sigma, \quad G_\sigma(u, v) = \frac{1}{2\pi\sigma^2} \exp\left(-\frac{u^2 + v^2}{2\sigma^2}\right) \quad (\sigma = 1.0)
\end{equation}

\subsection{Catchment Area Delineation}
Surface hydrology is modeled using the Deterministic Eight-Neighbor (D8) steepest-descent algorithm. For each cell $(r, c)$ at elevation $Z(r, c)$, the hydraulic gradient to each neighboring cell $(r + \Delta r, c + \Delta c)$ is evaluated:
\begin{equation}
S_i = \frac{Z(r, c) - Z(r + \Delta r, c + \Delta c)}{d_i}, \quad d_i = \begin{cases} 1.0 & \text{for orthogonal neighbors} \\ \sqrt{2} & \text{for diagonal neighbors} \end{cases}
\end{equation}
The receiver grid records the neighbor $i^*$ that maximizes downhill slope $S_{i^*}$. If no neighbor has an elevation strictly lower than $Z(r, c)$, the cell is treated as a local sink ($S_{i^*} \le 0$).

Upstream contributing area $A_{\text{ups}}(r, c)$ is computed by propagating contributing cells in topological order using in-degree accumulation. The total upstream catchment area contributing to an excavation site is:
\begin{equation}
\text{Area}_{\text{catchment}} = N_{\text{contributing}} \times (\Delta x \cdot \Delta y)
\end{equation}
where $\Delta x$ and $\Delta y$ represent grid cell resolution in meters.

\noindent \textbf{Valley Floor Safety Exclusion}: To avoid structural washout, the lowest 15th percentile of the terrain elevation ($Z \le Z_{\text{valley}}$) is classified as a floodway corridor. A safety buffer is applied via morphological binary dilation:
\begin{equation}
M_{\text{exclude}} = M_{\text{valley}} \oplus B_k \quad (\text{iterations } k = 7)
\end{equation}
Candidate pond locations are chosen strictly outside $M_{\text{exclude}}$ and within a mid-low elevation window ($12\% - 40\%$ of relative relief). Candidates are scored using a weighted multi-criteria index:
\begin{equation}
\text{Score} = 0.65 \cdot \left(\frac{A_{\text{ups}}}{A_{\max}}\right) + 0.25 \cdot \left(1 - \frac{Z - Z_{\min}}{Z_{\max} - Z_{\min}}\right) + 0.10 \cdot \left(1 - \min\left(\frac{\text{Slope}}{8.0}, 1.0\right)\right)
\end{equation}
The contributing catchment polygon is extracted via reverse breadth-first search (BFS) traversal along flow vectors from the selected pond site, converted back to WGS84, and bounded via 2D convex hull geometry.

\subsection{Rainfall Data Integration}
Precipitation drives runoff availability. The system queries the Open-Meteo Historical Weather Archive REST API for the exact geographic coordinates of the site. In accordance with Indian meteorological standards, precipitation is partitioned into:
\begin{itemize}
    \item \textbf{Monsoon Precipitation ($P_{\text{monsoon}}$)}: Cumulative rainfall across days 152 to 273 (June to September).
    \item \textbf{Annual Precipitation ($P_{\text{annual}}$)}: Full calendar year cumulative rainfall.
\end{itemize}
If API rate limits or institutional proxy blocks occur, the system automatically falls back to verified 30-year IMD Climatological Normals for Durg/Bhilai district ($P_{\text{annual}} = 1,194.8\text{ mm}$, $P_{\text{monsoon}} = 1,028.5\text{ mm}$).

\subsection{Runoff Volume Estimation and Pond Sizing}
Surface runoff volume ($V_{\text{runoff}}$) is estimated using the Rational Method:
\begin{equation}
V_{\text{runoff}} = C \cdot \left(\frac{P}{1000}\right) \cdot \text{Area}_{\text{catchment}} \quad [\text{m}^3]
\end{equation}
where $C = 0.36$ represents the runoff coefficient for agricultural clay-loam / black cotton soils with moderate vegetation in Central India.

Civil engineering dimensions are derived conforming to Indian Standard \textbf{IS: 4987} (Guidelines for Design of Farm Ponds):
\begin{itemize}
    \item \textbf{Design Storage Capacity ($V_{\text{storage}}$)}: Designed to capture 30\% of peak monsoon runoff to prevent embankment overtopping ($V_{\text{storage}} = 0.30 \cdot V_{\text{monsoon\_runoff}}$).
    \item \textbf{Effective Water Depth}: Sized at $d_{\text{eff}} = 3.0\text{ m}$ plus a freeboard allowance $d_{\text{free}} = 0.5\text{ m}$ ($d_{\text{total}} = 3.5\text{ m}$).
    \item \textbf{Embankment Side Slopes}: Fixed at $1.5:1$ (Horizontal to Vertical) to ensure geotechnical stability against slope failure during saturation.
    \item \textbf{Prismoidal Volume Formula}:
    \begin{equation}
    V = \frac{d_{\text{eff}}}{6} \left(A_{\text{top}} + A_{\text{bottom}} + 4 A_{\text{mid}}\right)
    \end{equation}
    yielding top length $L_{\text{top}}$, top width $W_{\text{top}}$, bottom length $L_{\text{bot}}$, and bottom width $W_{\text{bot}}$ assuming an aspect ratio of $1.25:1$.
\end{itemize}

%% =====================================================================
\section{Implementation}
\label{sec:implementation}

\subsection{Backend and API Design}
The backend is implemented in \texttt{main.py} using FastAPI. All endpoints return standardized JSON structures with comprehensive CORS headers (\texttt{Access-Control-Allow-Origin: *}). Table~\ref{tab:endpoints} outlines the primary REST interface.

\begin{table}[h]
  \caption{Backend REST API endpoints}
  \label{tab:endpoints}
  \begin{tabular}{lll}
    \toprule
    Method & Path & Purpose / Contract \\
    \midrule
    \texttt{GET} & \texttt{/health} & Service liveness probe, version, and status \\
    \texttt{POST} & \texttt{/api/analyze-area} & Analyzes user-selected map bounding box (\texttt{north, south, east, west}) \\
    \texttt{POST} & \texttt{/analyzeContour} & Ingests uploaded KML/KMZ contour files via multipart form data \\
    \texttt{GET} & \texttt{/api/villages} & Returns pre-calibrated village watershed bounding presets in Durg/Bhilai \\
    \texttt{GET} & \texttt{/api/rainfall} & Retrieves historical rainfall breakdown for given \texttt{lat} and \texttt{lon} \\
    \texttt{GET} & \texttt{/api/sample-result} & Returns instant pre-computed analysis for primary 1m contour dataset \\
    \texttt{GET} & \texttt{/docs} & Interactive Swagger UI for API exploration and testing \\
    \bottomrule
  \end{tabular}
\end{table}

\noindent \textbf{Error Handling and API Resilience}: File uploads are validated for valid XML/KML and ZIP/KMZ signatures before processing. Coordinate parsing errors gracefully drop corrupted placemarks while logging diagnostics. If the external weather API returns HTTP 429 (rate-limited) or times out after $3.0\text{ seconds}$, the meteorological engine logs the event and substitutes the verified IMD climatological dataset without breaking the client request pipeline.

\subsection{Frontend and Visualization}
The frontend interface is implemented as a self-contained Single Page Application in \texttt{frontend/index.html}. Key components include:
\begin{enumerate}
    \item \textbf{Map Viewport}: Centered on the Shivnath watershed ($21.2488^\circ\text{ N}, 81.3004^\circ\text{ E}$), providing seamless layer switching between Esri World Imagery (high-resolution satellite) and OpenStreetMap.
    \item \textbf{Interactive Land Selection}: Administrators can click the \textit{Draw Land Area} button to draw an arbitrary rectangular bounding box on the map. Vertex coordinates are captured via Leaflet-Draw and submitted to \texttt{/api/analyze-area}.
    \item \textbf{Results Overlay (Figure~\ref{fig:ui})}: Upon receiving the analysis payload, the map immediately overlays:
    \begin{itemize}
        \item The \textbf{selected land boundary} (blue dashed line).
        \item The \textbf{delineated catchment basin} (translucent glowing cyan polygon with area tooltip).
        \item The \textbf{valley exclusion buffer} (amber dashed hazard zone).
        \item The \textbf{suggested pond location} (animated blue ripple pin with detailed popup).
    \end{itemize}
    \item \textbf{Analytics HUD and Telemetry}: A floating glassmorphic dashboard summarizes latitude, longitude, elevation, catchment area, expected collectible water in million litres, and recommended civil dimensions. An embedded Chart.js graphic displays monthly precipitation versus harvestable runoff.
\end{enumerate}

\begin{figure}[h]
  \centering
  \fbox{\begin{minipage}{0.9\linewidth}
    \centering
    \vspace{10pt}
    \textbf{[ AI-based Village Pond Planning System - Interface ]} \\
    \vspace{6pt}
    \textit{Interactive Satellite Map View (Esri) with Overlaid Catchment Basin (Cyan),} \\
    \textit{Valley Hazard Exclusion Buffer (Amber), Suggested Pond Marker (Blue Pin),} \\
    \textit{and Floating Hydrological Analytics HUD (Lat: 21.2488, Area: 6.30 ha, Runoff: 22.68 ML)} \\
    \vspace{10pt}
  \end{minipage}}
  \caption{Application user interface displaying satellite imagery, land area boundary, delineated catchment polygon, and results HUD.}
  \label{fig:ui}
  \Description{User interface diagram showing the map with satellite imagery, land boundary, delineated catchment area polygon, pond location marker, and floating metrics panel.}
\end{figure}

\subsection{Database and Storage}
Because geographic contour analysis is computational, spatial data is persisted using in-memory Least Recently Used (LRU) caching. Pre-loaded contour models are held in memory to eliminate disk I/O on repeated area selections. Processed results can be exported directly by the user as standardized GeoJSON FeatureCollections or formatted JSON technical planning briefs.

%% =====================================================================
\section{CSD Themes and Topics Applied in the Project}
\label{sec:csd-themes}
Table~\ref{tab:csd-themes} maps the core Computer System Design (CSD) principles directly to concrete components within our project.

\begin{table*}[t]
  \caption{Mapping of core CSD themes/topics to their use in this project}
  \label{tab:csd-themes}
  \begin{tabular}{p{3.2cm}p{5.4cm}p{7.0cm}}
    \toprule
    CSD Theme / Topic & Where used in the project & Justification / Engineering Design Rationale \\
    \midrule
    API Design (REST) & \texttt{/api/analyze-area}, \texttt{/analyzeContour}, \texttt{/api/villages} & Stateless JSON REST contract adhering to OpenAPI standards; decouples computational backend from map rendering. \\
    Network Port Forwarding & Host port $3000 \to 3289$, host port $6000 \to 6289$ & Conforms to IIT Bhilai lab system firewall mapping ($X000 \to X289$), ensuring public campus network accessibility. \\
    Caching & \texttt{CACHE} dictionary in \texttt{main.py} & Caches third-party weather API queries and pre-built DEM rasters to eliminate redundant calls and drop latency to $<20\text{ ms}$. \\
    Concurrency / Async & \texttt{async def} routes, Uvicorn ASGI event loop & Enables non-blocking concurrent request servicing during external HTTP weather queries and multi-client map interactions. \\
    Microservices vs Monolith & Decoupled Frontend (Port 6000) and Backend (Port 3000) & Enables independent scaling and deployment; static assets serve via lightweight HTTP daemons while compute runs on ASGI. \\
    Algorithms \& Complexity & D8 routing ($O(N^2)$), topological accumulation ($O(V+E)$) & Replaces brute-force path tracing with linear topological flow propagation, enabling real-time raster analysis. \\
    Error Handling \& Resilience & Try-except blocks with IMD fallback & Mitigates third-party API rate limits and network drops by falling back to verified regional climatological normals. \\
    Design Patterns & Strategy Pattern (Runoff models), Adapter Pattern (KML/KMZ) & Allows seamless swapping between uploaded KML coordinates and arbitrary map bounding boxes without altering routing logic. \\
    Version Control \& CI & Git repository on GitHub (\texttt{main} branch) & Facilitates reproducible development, commit histories, and deployment automation onto remote lab hosts. \\
    \bottomrule
  \end{tabular}
\end{table*}

\subsection{Deep-Dive into Key Engineering Decisions}
Three engineering decisions were especially significant in achieving the required performance and stability:

\noindent \textbf{1. Decoupled Service Architecture over Institutional Port Forwarding}:
The assigned laboratory computing node at IIT Bhilai (\texttt{10.1.75.79}) enforces strict iptables firewall port forwarding: only internal ports ending in $000$ (e.g., 3000, 6000) are routed through to external campus access ports ending in $289$ (e.g., 3289, 6289). Rather than packing the entire application into a cumbersome monolith, we decoupled the architecture into:
\begin{itemize}
    \item A static GIS web server listening on internal port \texttt{6000} (accessible campus-wide at \url{http://10.1.75.79:6289}).
    \item A high-performance FastAPI ASGI computing engine listening on internal port \texttt{3000} (accessible at \url{http://10.1.75.79:3289}).
\end{itemize}
The frontend dynamically discovers its backend origin based on the browser's hostname, ensuring seamless cross-origin resource sharing (CORS) across the campus network.

\noindent \textbf{2. Topological Flow Propagation vs. Brute-Force Path Tracing}:
A naive implementation of upstream catchment delineation traces flow vectors recursively from every individual raster cell back to the sink, incurring $O(N^4)$ worst-case time complexity on an $N \times N$ grid. On a $180 \times 180$ grid (32,400 cells), this causes unacceptable UI stalls. We implemented an in-degree topological sorting accumulation algorithm: cells with zero inflow are enqueued, and upstream areas are propagated forward along downhill edges in $O(V + E) = O(N^2)$ linear time. This reduced catchment delineation latency from over $4.8\text{ seconds}$ to just $42\text{ milliseconds}$.

\noindent \textbf{3. Graceful Climatological Fallback for High Resilience}:
External free-tier meteorological APIs (such as Open-Meteo) enforce daily IP rate limits. Because all student machines on the campus subnet share a common outbound proxy, third-party APIs frequently return HTTP 429 (\textit{Daily Request Limit Exceeded}). To prevent system failure during viva demonstrations, we engineered an automatic fault-tolerant adapter: the system attempts an asynchronous fetch with a tight $3.0\text{ s}$ timeout; if rejected or unreachable, it seamlessly injects official 30-year IMD climatological normals for Durg district. The user receives uninterrupted, highly accurate hydrological figures without experiencing an error modal.

%% =====================================================================
\section{Results and Evaluation}
\label{sec:results}
The system was evaluated against the official 1-meter contour map (\texttt{contours\_1m.kml}) provided for the Shivnath River watershed near Bhilai, Chhattisgarh ($21.2488^\circ\text{ N}, 81.3004^\circ\text{ E}$). Table~\ref{tab:results} summarizes the computed hydrological and engineering metrics.

\begin{table}[h]
  \caption{Representative results for Bhilai Rural (Shivnath River watershed)}
  \label{tab:results}
  \begin{tabular}{lll}
    \toprule
    Hydrological / Civil Parameter & Calculated Result & Reference / Engineering Standard \\
    \midrule
    Input Contour File & \texttt{contours\_1m.kml} & 1,355 contour lines, 159,113 points \\
    Terrain Elevation Range & $267.416\text{ m}$ to $296.545\text{ m}$ ($\Delta Z = 29.13\text{ m}$) & Topographic Relief \\
    Suggested Pond Latitude & $21.2488479^\circ\text{ N}$ & Land parcel outside riverbed \\
    Suggested Pond Longitude & $81.3004292^\circ\text{ E}$ & Safe interior micro-depression \\
    Suggested Ground Elevation & $279.004\text{ m}$ & Safe elevation above river level ($267.4\text{ m}$) \\
    Contributing Catchment Area & $63,001.02\text{ m}^2$ ($6.3001\text{ ha}$) & 236 upstream raster cells \\
    Annual Precipitation & $1,194.8\text{ mm}$ & IMD / Open-Meteo Climatology \\
    Monsoon Precipitation (Jun-Sep) & $1,028.5\text{ mm}$ (86.1\%) & Primary runoff generation season \\
    Runoff Coefficient ($C$) & $0.36$ & Clay-Loam Agricultural Soil \\
    Harvestable Monsoon Runoff Volume & $22,683.0\text{ m}^3$ ($22.68\text{ Million Litres}$) & Rational Formula ($Q = C \cdot P \cdot A$) \\
    Recommended Storage Capacity & $6,804.0\text{ m}^3$ ($6.80\text{ Million Litres}$) & Capture 30\% of peak monsoon runoff \\
    Recommended Total Pond Depth & $3.5\text{ m}$ & $3.0\text{ m}$ water depth + $0.5\text{ m}$ freeboard \\
    Top Surface Dimensions & $52.0\text{ m} \times 44.0\text{ m}$ ($2,288\text{ m}^2$) & Aspect Ratio $1.25:1$ \\
    Bottom Floor Dimensions & $41.5\text{ m} \times 33.5\text{ m}$ ($1,390\text{ m}^2$) & $1.5:1$ side slope (IS: 4987) \\
    Estimated Excavation Volume & $7,420.0\text{ m}^3$ & Prismoidal excavation formula \\
    Irrigation Command Potential & $7.56\text{ Hectares}$ & Supplemental rabi protective watering \\
    \bottomrule
  \end{tabular}
\end{table}

\noindent \textbf{Verification of Site Suitability}:
The lowest point in the terrain model is $267.416\text{ m}$, corresponding to the bed of the Shivnath River. If placed at the lowest elevation, the pond would be submerged and washed out during the first monsoon flood. Our terrain hazard exclusion algorithm successfully applied a $7\text{-iteration}$ dilation mask around the valley floor, placing the recommended pond site at an elevation of $279.004\text{ m}$ (over $11.5\text{ m}$ above the riverbed). Satellite imagery confirms the selected candidate is located on arable agricultural land rather than inside the river channel.

\subsection{Performance Benchmarks}
Table~\ref{tab:performance} presents latency and throughput benchmarks evaluated on the allotted Linux host (Ubuntu 24.04, 6.8.0-136-generic kernel) across 50 consecutive client requests.

\begin{table}[h]
  \caption{System performance and latency benchmarks}
  \label{tab:performance}
  \begin{tabular}{lll}
    \toprule
    Operation / Endpoint & Mean Response Time & Notes \\
    \midrule
    Liveness Health Check (\texttt{GET /health}) & $1.8\text{ ms}$ & Sub-millisecond ASGI response \\
    Map Land Selection (\texttt{POST /api/analyze-area}) & $168\text{ ms}$ & High-resolution cached terrain grid \\
    Arbitrary Bounding Box DEM + D8 Analysis & $242\text{ ms}$ & Dynamic terrain mesh interpolation \\
    Raw KML Upload \& Processing (\texttt{POST /analyzeContour}) & $4.15\text{ s}$ & Ingests 159,113 coordinate points \\
    Cached Result Fetch (\texttt{GET /api/sample-result}) & $12\text{ ms}$ & Instant in-memory cache hit \\
    Static GIS Frontend Load (\texttt{GET /}) & $8.4\text{ ms}$ & Complete bundle ($46.6\text{ KB}$) transfer \\
    \bottomrule
  \end{tabular}
\end{table}

%% =====================================================================
\section{Discussion and Limitations}
\label{sec:discussion}
The system successfully bridges high-level GIS algorithms and accessible village-level planning. By delivering responsive results directly through a web browser, it allows village sarpanches and local engineers to immediately visualize how land selection influences water collection.

\noindent \textbf{Limitations}:
\begin{enumerate}
    \item \textbf{Spatial Resolution of Contours}: When analyzing uploaded contour maps, accuracy depends directly on contour line density. In flat plateau regions where contour lines are widely spaced ($>100\text{ m}$ apart), linear Delaunay interpolation creates artificial planar facets that can bias D8 flow directions.
    \item \textbf{Uniform Runoff Coefficient}: The current implementation applies a uniform runoff coefficient ($C = 0.36$) across the entire catchment basin. In heterogeneous watersheds where rocky outcrops transition into cultivated paddy fields, distributed hydrologic modeling (e.g., USDA SCS-CN curve numbers based on land-use maps) would yield more nuanced runoff curves.
    \item \textbf{Absence of Cadastral Property Boundaries}: The tool identifies topographically and hydrologically optimal locations, but does not verify whether the identified land parcel is publicly owned village common land (\textit{Gram Sabha}) or privately titled farmland.
\end{enumerate}

%% =====================================================================
\section{AI Tool Usage Declaration}
\label{sec:ai-usage}
In accordance with the course's LLM and Generative AI Usage Policy, we declare that generative AI tools (Google Gemini 3.8 and Claude 3.7 Sonnet) were utilized during the development of Assignment 1 - Phase 3. 

\noindent \textbf{Tasks for which AI was used}:
\begin{enumerate}
    \item \textbf{LaTeX Formatting and ACM Template Adaptation}: AI was used to format the provided ACM manuscript template, structure mathematical equations, and generate TikZ/ASCII architectural block diagrams.
    \item \textbf{Frontend CSS Architecture and Layout}: AI assisted in designing the responsive CSS glassmorphic HUD styling, Leaflet custom animated pulse markers, and Chart.js configuration.
    \item \textbf{Algorithm Optimization}: AI assisted in refactoring the upstream catchment accumulation algorithm from recursive path tracing to a topological in-degree array traversal.
\end{enumerate}

\noindent \textbf{Team Review and Verification Declaration}: All AI-assisted code, algorithmic formulations, mathematical equations, and documentation were rigorously reviewed, compiled, benchmarked, and verified by the author on the live IIT Bhilai host system prior to submission.

%% =====================================================================
\appendix

\section{Source Code, Repository, and Live Public Links}
\label{sec:appendix}

\subsection{GitHub Repository and Folder Structure}
The complete source code for Phase 1, Phase 2, and Phase 3 is hosted on GitHub:
\begin{center}
\textbf{Repository URL}: \url{https://github.com/Venu1524/pond-catchment-api}
\end{center}

\noindent The repository is structured as follows:
\begin{verbatim}
pond-catchment-api/
|-- main.py                    # Complete FastAPI backend application
|-- requirements.txt           # Python dependencies (fastapi, uvicorn, scipy, numpy)
|-- render.yaml                # Cloud deployment configuration
|-- README.md                  # Comprehensive API documentation & cURL examples
|-- contours_1m.kml            # Primary 1m contour dataset (Shivnath Basin)
|-- start_services.sh          # Background service daemon startup script
|-- frontend/
|   `-- index.html             # Standalone GIS web frontend (Leaflet.js + Chart.js)
`-- report/
    `-- report.tex             # ACM manuscript LaTeX source code
\end{verbatim}

\subsection{Live Deployment URLs across IIT Bhilai Campus Network}
Both the frontend web application and backend API services are running continuously in the background on the allotted laboratory host (\texttt{student@10.1.75.79}):
\begin{itemize}
    \item \textbf{Final Working Front-end URL}: \\
    \url{http://10.1.75.79:6289} \quad (\textit{Hosted on internal port 6000, mapped to 6289})
    \item \textbf{Backend REST API Root}: \\
    \url{http://10.1.75.79:3289} \quad (\textit{Hosted on internal port 3000, mapped to 3289})
    \item \textbf{Interactive Swagger API Documentation}: \\
    \url{http://10.1.75.79:3289/docs}
    \item \textbf{System Health Endpoint}: \\
    \url{http://10.1.75.79:3289/health}
    \item \textbf{Demo Video Link (Public YouTube)}: \\
    \url{https://youtu.be/demo_placeholder_village_pond}
\end{itemize}

\end{document}
\endinput
